

% First define the torus of size $2^n$ and in dimension $d$,
% and introduce the Ising model on it.
% Also define the discrete GFF with a magnetic field on it.

\section{Reflection positivity}

\begin{definition}[Torus]
  \label{def:torus}
    For $n\geq 1$ and $d\geq 1$, we define the torus of circumference $n$ in dimension $d$ as the graph $\T_n^d$ with vertex set
    \(
        \T_n^d := (\Z/n\Z)^d
    \)
    and edge set
    \[
        E(\T_n^d) := \{\{x,y\} : x,y\in \T_n^d, \|x-y\|_1 = 1\}.
    \]
    We occasionally write $\T$ for $\T_n^d$ when the parameters are clear from the context.
\end{definition}

% Now state the chessboard estimate, for the Gaussian free fied.

\begin{theorem}[Chessboard estimate]
  \label{thm:chessboard}
  Fix $d\in\Z_{\geq 1}$ and $n=2^k$ for some $k\in\Z_{\geq 1}$.
  Also fix some $\beta\geq 0$.
  Then for any family of real functions $f_x:\R\to\R$ indexed by $x\in \T$, we have
  \[
    \left|\int_{\R^\T}
    e^{-\frac\beta2 \|\nabla\phi\|_2^2} \prod_{x\in \T} f_x(\phi_x) \diff\phi\right|
    \leq
    \prod_{x\in \T}\left( \int_{\R^\T}
    e^{-\frac\beta2 \|\nabla\phi\|_2^2} \prod_{y\in \T} f_x(\phi_y) \diff\phi\right)^{1/|\T|},
  \]
  and in particular the right side is bounded by
  \[
  \max_{x\in \T}
    \int_{\R^\T}
    e^{-\frac\beta2 \|\nabla\phi\|_2^2} \prod_{y\in \T} f_x(\phi_y) \diff\phi.
  \]
\end{theorem}

We apply the theorem in the following way:
let $\tau\in\R^\T$ and $A\subset\R$, and set $f_x(a) := \true{a-\tau_x\in A}$ for each $x\in \T$.
Then the chessboard estimate gives 
\[
    \int_{\R^\T}
    e^{-\frac\beta2 \|\nabla\phi\|_2^2} \true{\phi-\tau\in A^\T} \diff\phi
    \leq
    \int_{\R^\T}
    e^{-\frac\beta2 \|\nabla\phi\|_2^2} \true{\phi\in A^\T} \diff\phi.
\]
Now suppose that we set $A=[-1,1]\setminus(-1+\epsilon,1-\epsilon)$,
multiply both sides by $\epsilon^{-|\T|}$, and take the limit $\epsilon\to 0$.
Then the right side converges to
\[
    \sum_{\sigma\in\{\pm 1\}^\T} e^{-\frac\beta2 \|\nabla\sigma\|_2^2}
    =
    e^{\beta |E(\T)|}
    Z_{\T,\beta}^{\operatorname{Ising}},
\]
while the left side converges to
\[
    \sum_{\sigma\in\{\pm 1\}^\T} e^{-\frac\beta2 \|\nabla(\sigma+\tau)\|_2^2}=
    e^{\beta |E(\T)|}
    Z_{\T,\beta}^{\operatorname{Ising}}\cdot
    \langle
    e^{\frac\beta2 (\|\nabla\sigma\|_2^2-\|\nabla (\sigma+\tau)\|_2^2)}
    \rangle_{\T,\beta}^{\operatorname{Ising}}.
\]
In particular, the expectation on the right is bounded by $1$, which leads to
\[
    \langle
    e^{\beta (\nabla\sigma,\nabla\tau)}
    \rangle_{\T,\beta}^{\operatorname{Ising}}
    \leq e^{\frac\beta2 \|\nabla\tau\|_2^2}.
\]
Replacing $\tau$ by $\epsilon\tau$, expanding to second order in $\epsilon$, and sending $\epsilon\to0$,
yields
\[  
    \frac{\beta^2}2
    \langle (\nabla\sigma,\nabla\tau)^2\rangle \leq \frac\beta2 \|\nabla\tau\|_2^2
    \qquad
    \text{or}
    \qquad
    \langle (\nabla\sigma,\nabla\tau)^2\rangle \leq \frac1\beta \|\nabla\tau\|_2^2
\]

Now define the Laplacian $\Delta$ on $\T$ by
\[
    (\Delta f)_x := \sum_{y\sim x} (f_y - f_x).
\]
Then for any $f,g\in\R^\T$, we have
\begin{equation*}
    (\nabla f,\nabla g) = \sum_{\{x,y\}\in E(\T)} (f_x-f_y)(g_x-g_y) = -\sum_{x\in \T} f_x (\Delta g)_x = -(f,\Delta g).
\end{equation*}
Then the previous inequality can be rewritten as
\begin{equation}
    \label{eq:Gaussian_domination_real}
    \langle
    (\sigma,\Delta\tau)^2
    \rangle
    \leq \frac1\beta(\tau,-\Delta\tau).
\end{equation}

Notice that $-\Delta:\R^\T\to\R^\T$ is a linear operator commuting with shifts
of the torus.
This suggests that it is diagonal in the Fourier basis.

\color{blue}
Recall that $\T=(\Z/n\Z)^d$.
Let \(\operatorname{Char}(\T)\) be the set of \emph{characters}
of $\T$, that is, the set of group homomorphisms $(\T,+)\to\C$.
Since $\T$ is a finite abelian group, \(\operatorname{Char}(\T)\) may be understood as follows.
\begin{itemize}
    \item Since $nx=0\in\T$ for any $x\in\T$, we have $\chi^n\equiv 1$ for any $\chi\in\operatorname{Char}(\T)$.
    \item If we assign an $n$-th root of unity to each basis element $(e_i)_i\in\T$, then we get a character by extending the assignment linearly.
\end{itemize}
This gives a description of all characters.
More precisely,
writing $\T^*:=(2\pi/n)(\Z/n\Z)^d$ for the \emph{dual torus}, we can identify $\operatorname{Char}(\T)$ with $\T^*$ by sending $k\in\T^*$ to the character $\chi_k$ defined by $\chi_k(x):=e^{i(x,k)}$.

Now $\ell^2(\T)$ is a complex Hilbert space with inner product
\[
    (f,g):=\sum_{x\in \T} \bar g_x f_x.
\]
Moreover, $\operatorname{Char}(\T)$ forms an orthogonal basis of $\ell^2(\T)$, with $\|\chi\|_2^2=|\T|$ for each character $\chi$.
This means that for any $f\in\ell^2(\T)$, we have
\[
    f = \sum_{\chi}\frac{1}{\|\chi\|_2^2}(\chi,f)\chi = \frac1{|\T|}\sum_{\chi}(\chi,f)\chi.
\]
The vectors in the basis $\operatorname{Char}(\T)$ are eigenvectors of translations of the torus.
For this reason, it is also called the \emph{Fourier basis}.
For any $f\in\ell^2(\T)$ and $k\in\T^*$, we write $\hat f(k):=(\chi_k,f)$ for the \emph{Fourier coefficient} of $f$ at $k$.
The formula above may then be rewritten as
\[
    f = \frac1{|\T|}\sum_{k\in\T^*} \hat f(k) \chi_k.
\]
This is an orthogonal expansion of $f$ (with $\|\chi_k\|_2^2=|\T|$), which implies Parseval's identity:
\[
    \|f\|_2^2 = \frac1{|\T|}\sum_{k\in\T^*} |\hat f(k)|^2.
\]
Since the shift operators on $\T$ are diagonal in the Fourier basis,
all convolution operators on $\ell^2(\T)$ are also diagonal in the Fourier basis.
This includes the Laplacian $\Delta$.
\color{black}

By a direct calculation, one can check that
\[
    -\Delta\chi_k = \lambda_k\chi_k
    \qquad\text{with}\qquad
    \lambda_k = 2\sum_{i=1}^d (1-\cos k_i).
\]


Equation~\eqref{eq:Gaussian_domination_real} only makes sense for real-valued $\tau$,
but we would like to set $\tau=\chi_k$.
By applying Equation~\eqref{eq:Gaussian_domination_real} explicitly for the
real and imaginary part of $\chi_k$, we obtain
\begin{equation}
    \label{eq:Gaussian_domination_Fourier}
    \lambda_k^2\langle|(\sigma,\chi_k)|^2\rangle \leq \frac{\lambda_k}\beta |\T|
    \qquad\text{or}\qquad
    \langle|(\sigma,\chi_k)|^2\rangle \leq \frac{\T}{\beta\lambda_k}.
\end{equation}
Now for fixed $\sigma\in\{\pm 1\}^\T\subset \R^\T$,
we may interpret $(\hat\sigma_k)_k:=((\sigma,\chi_k))_k$ as the Fourier
transform of $\sigma$.

\subsection{Alternative proof of spontaneous magnetization for the Ising model in dimension $d\geq 3$}

Parseval's identity yields
\[
    \sum_k |\hat\sigma_k|^2 =  |\T|(\sigma,\sigma)| = |\T|^2.
\]
By Equation~\eqref{eq:Gaussian_domination_Fourier},
we may now \emph{lower bound} $\langle|\hat\sigma_0|\rangle$
by
\begin{equation}
    \frac{\langle|\hat\sigma_0|^2\rangle}{|\T|^2}\geq 1 - \frac1{\beta|\T|}\sum_{k\in\Lambda^*\setminus\{0\}}\frac{1}{\lambda_k}.
\end{equation}
Now we may interpret the right-hand side as a Riemann sum,
to obtain the lower bound 
\[
    \frac1{|\T|}\sum_{k\in\Lambda^*\setminus\{0\}}\frac{1}{\lambda_k}
    \geq
    (2\pi)^{-d}\int_{[-\pi,\pi]^d} \frac1{2\sum_{i=1}^d (1-\cos t_i)}\diff t=:I_d.
\]
By expanding the cosine to second order, one can check that $I_2=\infty$, but $I_d<\infty$ for $d\geq 3$.
Thus, we get
\[
    \frac{\langle|\hat\sigma_0|^2\rangle}{|\T|^2}
    \geq
    1-\frac{I_d}{\beta}>0
\]
for sufficiently large $\beta$.
This is great news,
since
\[
    \frac{\langle|\hat\sigma_0|^2\rangle}{|\T|^2}
    =
    \frac{1}{|\T|^2}\sum_{x,y\in \T} \langle\sigma_x\sigma_y\rangle,
\]
that is, it measures the two-point correlation function between two randomly chosen vertices of the torus.
It is easy to see that
\[
    \langle\sigma_x\rangle_{\Z^d,\beta}^+ \geq
    \sqrt{\frac{\langle|\hat\sigma_0|^2\rangle}{|\T|^2}}
    \geq  \sqrt{1-I_d/\beta}>0,
\]
which yields another proof of the phase transition for the Ising model on $\Z^d$ for $d\geq 3$.

\subsection{Upper bound on the two-point function at the critical point}

Let $S(x):=\langle\sigma_0\sigma_x\rangle$.
Its Fourier transform satisfies
\[
\hat S(k)=
|\sum_{x\in\T}\chi_k(x+y)\langle\sigma_0\sigma_x\rangle|
\leq
\langle|\sum_{x\in\T}\sigma_0\sigma_x\chi_k(x+y)|\rangle
=
\frac1{|\T|}\langle|\hat\sigma_k|^2\rangle
\leq
\frac{1}{\beta\lambda_k}.
\]
\todo{Explain why Fourier transform is non-negative}
We would like to turn this into a pointwise bound on $S(x)$.
We could express $S(x)$ in terms of its Fourier transform,
but this does not give a good bound.








% Now use the chessboard esimate to give the Gaussian domination bound